% Original: PDF strana 10, donji slajd; drugi deo.
\slajdid{03-s020b}
\begin{frame}[t,label=03-s020b]{Prenos dozvole signalom EO}
\setlength{\parskip}{0pt}
\fontsize{11}{13}\selectfont
% element: 03-s020b:e04
Na prvi pogled bismo umesto EO mogli koristiti invertovani GS.
Međutim, dozvola mora da se prenese \textbf{od višeg ka nižem prioritetu},
bez obzira na stanje pojedinačnog kodera.
\par\vspace{2mm}
\begin{columns}[T,onlytextwidth]
\begin{column}{.53\textwidth}
\textcolor{DEisticanje}{\textbf{Zašto invertovani GS nije dovoljan?}}
\par\vspace{1.5mm}
Uslovljeni GS može biti nula zato što je koderu \textbf{zabranjen rad}.
Njegova inverzija tada bi dozvolila nižem koderu da radi, \textbf{a ne sme}.
\par\vspace{1.5mm}
GS možemo ostaviti i neuslovljen: to ne smeta drugom nivou kodera prioriteta.
Ipak, ni on ne može zameniti EO. To što ova grupa nema šta da koduje
ne znači da \textbf{grupe višeg prioriteta nisu imale zahtev}.
\end{column}
\begin{column}{.44\textwidth}
\textcolor{DEisticanje}{\textbf{Suština izlaza EO}}
\par\vspace{1.5mm}
\textbf{Dozvoljeno mi je; nemam šta da kodujem:}
dozvoljavam kodovanje nižem koderu.
\par\vspace{2mm}
\textbf{Dozvoljeno mi je; imam šta da kodujem:}
zabranjujem kodovanje nižim koderima.
\par\vspace{2mm}
\textbf{Nije mi dozvoljeno da kodujem:}
ne dozvoljavam ni nižim koderima da koduju.
\end{column}
\end{columns}
\end{frame}
